
            \documentclass{article}
            \usepackage[utf8]{inputenc}
            \usepackage{hyperref}
            \usepackage{verbatim}

            \begin{document}

            \section{Configure a dataset}

Before creating a dataset, the user has the possibility to 
specify the preferences of the calibrations to associate, and configure the 
workflow parameters to select the reduction strategies. 

\subsection{Select the desired reduction target}
The "Reduction Target", or the "target task" is the end point of the reduction cascade. When specifying a target, EDPS will process
all and only the files needed to execute it. For example, if my target is "science", and the science files
need the bias files, EDPS will process only the biases that have been selected to process those science files; then it 
processes the science using the product of the bias reduction.
However, if my target is bias, then EDPS will process all and only the bias files, regardless they are not used by any science. 
In this case, EDPS does not process the science, as it has reached already the end reduction point (e.g., process all biases).
The "Target category" is a group of target that have similar purposes. For example, the target category "science", includes
all the tasks that deliver final scientific products, the target category "qc1calib" includes all and only the tasks that 
processes calibrations (e.g., bias, standard stars).

For scientific data reduction purposes, leave the default Target Category and Reduction Targets. 


\begin{verbatim}
:alt: select_target
:name: select_target
\end{verbatim}


\subsection{Select appropriate calibrations}
\begin{verbatim}
\end{verbatim}

\subsection{Configure workflow parameters}

This window allows to:
\begin{itemize}
\item Select the parameter set. A pre-determined list of workflow parameters and recipe parameters for a
\end{itemize}
  given use case. For the majority of the cases, the “science” parameter set can be used.
\begin{itemize}
\item Edit the workflow parameters. These are parameters that regulates the
\end{itemize}
  reduction strategy, e.g. whether to use a given calibration or not, or to trigger a certain reduction step. 
  Some changes can define a strategy that selects different files; in this case a new dataset is created.


\begin{verbatim}
:alt: select_wkf_parameters
:name: select_wkf_parameters
\end{verbatim}

            \end{document}
